%% kapa: exam=1 round=30 year=2019 period=1 form=A subject=economics subject_ko=경제학원론 seq=13 number=53
%% source: Q-Net 감정평가사 기출문제 (cst00309) artlSeq=5205205 file=extracted/1008/5205205/2019년 제30회 감정평가사 1차시험 문제지/2019년 제30회 감정평가사 1차시험 문제지/2019년도 제30회 감정평가사 1차 1교시 A형.hwp
%% conversion: kordoc 4.24.51 -> tools/gen1.py (auto); review=proofread-sonnet; flags=none
\begin{question}{53}
완전경쟁시장에서 개별기업은 U자형 평균비용곡선과 평균가변비용곡선을 가진다. 시장가격이 350일 때, 생산량 50 수준에서 한계비용은 350, 평균비용은 400, 평균가변비용은 200이다. 다음 중 옳은 것을 모두 고른 것은?
\begin{framedtext}
ㄱ. 평균비용곡선이 우상향하는 구간에 생산량 50이 존재한다. \par
ㄴ. 평균가변비용곡선이 우상향하는 구간에 생산량 50이 존재한다. \par
ㄷ. 생산량 50에서 음(-)의 이윤을 얻고 있다. \par
ㄹ. 개별기업은 단기에 조업을 중단해야 한다.
\end{framedtext}
\begin{choices}
  \item ㄱ, ㄴ
  \item ㄱ, ㄷ
  \item ㄱ, ㄹ
  \item ㄴ, ㄷ
  \item ㄴ, ㄹ
\end{choices}
\end{question}
\answer{4}
